\section{Institutional Background}
This section provides a brief institutional background on prescription drug litigation in the context of public tenders in Brazil.

First, we briefly introduce some relevant elements of the Brazil health litigation issue and its direct impacts on the public budget. We focus on the impacts of judicial decisions and administrative requests on the planning process of public tenders. Finally, we briefly describe the bidding process for prescription drugs, underlining the bidding negotiations’ difficulties resulting from court orders and administrative demands.

\subsection{Health: “Right of All and a Duty of the State”}
The 1988 Brazilian Federal Constitution created the Brazilian Unified Public Health System (SUS) to promote universal health coverage, t, which consists of a massive set of actions and programs jointly subsidized and implemented by the federal government, states, and municipalities. Although the SUS still has some issues and distortions, it has generally brought excellent public health results in Brazil \citep{castro2019, soares2019}.

One of the SUS’s main goals is to facilitate access to prescription and nonprescription drugs and other health items. However, this objective must meet public budget constraints, especially in emerging economies with a chronic fiscal deficit such as Brazil. In that respect, the SUS provides a list of procedures, medicines, and other health-related products that the government is committed to offering the population through its programs.

The SUS list works as a “social contract.” It is how Brazilian society deals with the trade-off between universal health coverage and public budget costs. Periodically, the SUS list is updated to keep up with technological changes in the health area and treatments of new and known diseases. The ultimate goal is to serve as many people as possible as long as the government maximizes new therapies’ cost-efficacy.

However, judges in Brazil tend to interpret the constitutional text literally and disregard costs in their analyses and decisions. Among many articles in the Constitution is a specific one (Article 196) that states: “Health is a right of all and a duty of the state and shall be guaranteed utilizing social and economic policies aimed at reducing the risk of illness and other hazards and at the universal and equal access to actions and services for its promotion, protection, and recovery.” This article is a general article that gives rise to a wide range of interpretations that bring significant distortions to the health system and proper public resources use.

Since the mid-2000s, the common understanding of judges has been that the government must provide all health items and procedures for the population at any time. This strict interpretation creates a detrimental scheme of incentives for different groups of agents. Individuals often sue the Brazilian state “[…] claiming that they have the right to receive the treatment they need or to be funded by the public health system” \citep{wang2015right}, whether or not the treatment is on the SUS list.

It is relatively easy and inexpensive to access the legal system in this context: individuals need only a prescription for the desired drugs and a private lawyer or public defender. In addition, the success rate of health litigation is very high: in the state of São Paulo, for example, approximately 85\% of first-instance claims prosper \citep{cnj2019judicializacao}, and the rates in superior courts are even higher.

The combination of the low cost of accessing the legal system and high success rates leads to strong incentives to obtain medicines through the courts. There has been a steady upward trend in judicial claims for health-related products in Brazil in recent years; first instances of this type of court order totaled almost 96,000 in 2017, increasing almost 130\% over 2008 \citep{cnj2019judicializacao}. In the same period, the growth in the total number of lawsuits in the lower courts was only approximately 50\%.

In the state of São Paulo, the growth pattern was even higher. An approximate increase of 913\% occurred between 2008 and 2017, increasing from 2,317 to 23,465 yearly lawsuits for health products \citep{cnj2019judicializacao}.

The lawsuits occur in many municipalities with wide dispersion throughout the state of São Paulo (see Figure~\ref{fig:distribution}). A higher concentration of cases occurs in the most populous municipalities in absolute terms.

\begin{figure}[ht]
  \centering
  \includegraphics[width=0.8\textwidth]{figure1.png}
  \caption{Distribution of Health Litigation Cases – Municipalities in the State of São Paulo (2008--2017)\\ \emph{Source: S-CODES (SES/SP) and CNJ/INSPER (2019)}}
  \label{fig:distribution}
\end{figure}

In addition to wide spatial dispersion, court orders consist of a massive variety of different items across time, including differences in dosages and drug presentations. Between 2009 and 2018, approximately 2,760 different items were ordered at least once a year on average.

The large and increasing number of successful judicial requests strongly affects the public health budget. The state government of São Paulo has a total annual budget of \$58~billion, of which approximately 10\% (\$5.9~billion) goes to funding the Public Health System in the state. In 2018, public spending only for litigated health items was nearly 5\% of the annual budget (US\$300~million) of the Department of Public Health of São Paulo state (SES/SP).

The SES/SP is responsible for managing public resources and implementing policies and programs to promote public health in the state of São Paulo. It consists of 99 decentralized public buyer units (PBUs) distributed throughout the state (see Figure~\ref{fig:pbu}).
  
\begin{figure}[ht]
  \centering
  \includegraphics[width=0.8\textwidth]{figure2.png}
  \caption{Public Buyer Units (SES/SP)}
  \label{fig:pbu}
\end{figure}

Most of the SES/SP’s annual budget goes to purchase common goods and services, especially prescription drugs and other health-related items, to support all public health programs in São Paulo state. PBUs are directly responsible for planning and making those purchases.

The purchase process is greatly affected when there is a court order to acquire a specific item. The court order obliges PBUs to buy quickly, with massive restrictions on the planning process and heavy sanctions against public officials if they do not comply with the judicial order.

\subsection{Planning under Pressure: Judicial Decisions and Administrative Requests as Restrictions for Planning}
As in many other countries, Brazilian law establishes as a general rule that all purchases, services, and works hired by the public administration are subject to a public tender. Federal Law 8,666/1993 institutes the general framework applicable to all public bids in the country, which must be observed by all three government branches.

According to Law 8,666/1993, a public purchase comprises three distinct, mandatory, and subsequent phases: (i) the internal phase (planning and publication of the notice), (ii) the external phase (negotiation between purchasing units and suppliers), and (iii) delivery of items.

The internal phase of an ordinary purchase consists of the public administration carrying out careful procurement preparation and planning. At this stage, PBUs first identify their needs and select what types of goods or services might be appropriate to meet those demands. Guided by the SES/SP headquarters, the purchases made by each PBU take into account local and regional demands. However, the most important requirement is that the SES/SP exclusively purchase items that appear on the SUS list.

The tender preparation commission then creates a purchase order defining the main parameters of the bidding process. These parameters consist of the number of items to be purchased and their specifications, the bidding schedule for all participants, the bidding procedure, the auctioneer in charge, quantities, reference prices, delivery addresses, and minimum bidder requirements for participation, payment method, and possible fines. All of these parameters, except reference prices, are brought together in a document called a public notice.

Choosing suitable bid parameters, such as quantities and reference prices, increases the chances of an efficient purchase. PBUs need enough time to accurately organize a proposal and take advantage of the government’s ability to buy items on a large scale at better prices. Nonetheless, in some situations, planning time is scarce.

When court decisions reach the government, they force PBUs to buy items in very adverse conditions. First, these healthcare-related court decisions are almost exclusively enforced through injunctions (99.94\% of all court decisions). The injunctions’ application makes the deadlines for planning purchases and delivery of items very tight (between 1 and 10 days) compared to the deadlines for ordinary purchases (from 30 to 180 days). Consequently, the internal phase of litigated purchases is accomplished, on average, in one-third of the time of that of ordinary purchases, which undermines the process of setting the essential tender parameters.

Additionally, court orders have specific features that bring difficult and unpredictable elements to the purchase planning process. For example, since PBUs generally do not buy items that are not on the SUS list, they have less experience planning the purchase of these items. Most litigated products (almost 75\%), mainly high-cost products and those intended for severe or rare diseases, such as cancer and amyotrophic lateral sclerosis (ALS), are not listed in the SUS.

However, the courts also order prescription drugs that are on the SUS list. Most court orders require the whole treatment of a specific disease for an individual: a set or package of different prescription drugs. Almost 65\% of claims consist of items on the SUS list ordered along with non-SUS items. Although some of these items are on the SUS list, they also have a material impact on the public budget and purchase planning: litigated purchases use resources that would have been used for other purposes or were outside the initially planned budget.

Other relevant reasons for litigation are “off-label” uses of SUS-list items (approximately 20\%) and “jumping the line” in public programs (nearly 11\%). Additionally, judges heed complementary justifications such as individuals’ insufficient financial resources and imminent risk of death without requiring detailed evidence of these conditions in their decisions. Only 4\% of judicial claims are due to lack of stock or inability to provide service. Thus, this situation indicates that the “litigation shock” is poorly correlated with a possible unobservable characteristic of PBUs, such as mismanagement.

Only in rare cases can the government execute a court order using existing stock since (i) planning is performed to meet the demands of existing programs, and (ii) it is challenging to maintain and manage strategic stocks due to drugs’ high degree of perishability and the massive variety of items. For these reasons, health litigation acts as an exogenous shock, a severe restriction to be addressed in the planning process. It is not possible to anticipate exactly where, when, what, and what quantities the SES/SP might have to purchase. PBUs have little control over planning under these conditions. Additionally, health litigation has very little to do with the quality of the public policies implemented.

There are also other costs to the government resulting from the judiciary monitoring and actively participating in nonprescription purchases. Penalties are extremely severe for public officials if they do not comply with a court order. This potential threat constitutes an additional restrictive element to tenders generated by court orders. The primary forms of punishment for noncompliance are (i) fines (sometimes reaching significant and disproportionate amounts); (ii) administrative, civil and criminal liability; and (iii) blocking and “hijacking” public funds.

Since 2009, the SES/SP has tried to mitigate the monitoring and punishment costs generated by court orders. The mechanism used is negotiating an item’s supply directly with an individual before a court order occurs. This procedure is known as an administrative request. Administrative requests are evaluated by a scientific committee that can judge whether a request is valid. This commission uses the scientific literature with a healthy level of evidence, using evidence-based medicine criteria and protocols recognized by the medical community.

One main difference between a judicial and an administrative request is that the latter undergoes a scientific examination and tends to represent a better use of the resource and the drug. As it is a more rigid and scientific process, it tends to be less sought after and to generate fewer purchasing processes. The purchases generated by administrative orders totaled 9,700 between 2009 and 2019, representing approximately 6.5\% of the total purchases from court orders.

Administrative requests bring another benefit to public administration. There are no penalties for the public officials involved in the event of a failure in the requested item’s purchase process. Nevertheless, there is no difference between a court order and an administrative request in terms of planning purchases. Like a court order, an administrative request generates a purchase order with immediate delivery and using budget resources not designated for this purpose.

In general, purchases of prescription and nonprescription drugs can be classified into three groups: (i) ordinary, (ii) administrative, and (iii) litigated. Table~1 summarizes the types of purchases and their characteristics.


\begin{figure}[ht]
  \centering
  \includegraphics[width=0.8\textwidth]{figure3.pdf}
  \caption{Types of Purchases}
\end{figure}

Health litigation and administrative requests significantly impair the public procurement process, forcing public administration to respond to these demands promptly and without proper planning. The purchase and negotiation process itself is severely hampered.

\subsection{Buying under Pressure: Urgent Purchases as Public Information}
The external phase covers the time lapse between the publication of the public notice and the contractual signature. This phase involves interaction between the government and firms through a previously chosen competitive tendering procedure and other parameters defined in the internal phase.

The main objective of a bidding process conducted by a PBU is to seek the best possible contract for the government, taking into account the parameters defined in the planning phase. The public official responsible for negotiating with suppliers cannot change any previously defined parameters such as quantities, delivery time, reference prices, and tendering procedures.

The way the internal phase routines are performed may substantially affect the bidding process results. Given that urgent purchases (litigated or administrative purchases) are planned under very restrictive conditions, they can make negotiations very difficult. Therefore, the expected outcomes of these urgent purchases tend to be much worse than the outcomes when purchases are made under ordinary planning conditions.

In addition, informing all that a tender is an urgent purchase can amplify these effects. It is mandatory by law to provide information in the public notice that the purchase originated from a court order or administrative request; this is public information. For example, participating suppliers know that the SES/SP and its officers responsible for the bidding may be punished if a negotiation for a litigated purchase is unsuccessful. Making this situation public information thus can be a relevant imbalance factor in the bargaining process between the government and firms.

Despite the differences in planning conditions, ordinary and urgent purchases are made through the state e-platform under the same operational conditions. Each PBU purchases in a decentralized way through the \textit{Bolsa Eletrônica de Compras} (BEC), the e-procurement platform of São Paulo state. Since 2007, it has been mandatory to use the BEC to purchase common goods and services in São Paulo state, including all 99 SES/SP units. The BEC figures of SES/SP buying are expressive. In 2018, approximately US\$1.7~billion were traded, and since 2009, the electronic platform has handled more than R\$7.4~billion in SES/SP negotiations (approximately 34\% of all state purchases).

Almost 55,000 item purchase offers were negotiated, comprehending approximately 6,350 different traded items. The BEC has an extensive catalog of standardized items and services that are described in great detail.

The SES/SP uses the same two types of competitive tendering procedures available at the BEC to make ordinary and urgent (litigated and administrative) purchases: (i) sealed-bid tendering (\textit{convite}) and (ii) multiround descending auctions (\textit{pregão}).

In sealed bids, firms send their proposals to the government by a specific date specified in the notice. At a later date, the proposals and participants become public information when the auctioneer “opens” the envelopes. The winning firm is the one among those with appropriate documentation that submitted the lowest bid. \textit{Convite} is allowed up to the purchase value limit of R\$176,000 (approximately US\$35,200).

\textit{Pregão} has no limit on the purchase value. This mode is a combination of a modified sealed-bid tender and reverse auction. In this case, PBUs rank the qualified proposals only under the conditions set out in the sealed-bid phase’s notice. Then, the auctioneer publicly reveals all valid proposals, keeping firms’ identities anonymous.

Next, the descending auction begins: for 20 minutes, each qualified firm submits its bids, knowing the current lowest valid bid. If there is a valid bid between 16 and 20 minutes, the auction will be extended for another 4 minutes. It ends only if 4 minutes passed with no valid bid. The final criterion for winning the tender is presenting the lowest bid price, which must be lower than the reference price.

One of the main differences between the \textit{convite} and \textit{pregão} modes is that the latter allows for a negotiation phase after the reverse auction during which companies and the government can negotiate the lowest final price previously obtained. On the other hand, since \textit{convite} has a single phase, it tends to be a more straightforward procedure to perform and monitor.

Planning and executing tenders consist of a very costly public administration process that demands relevant financial and human resources. An acquisition made by \textit{pregão} or \textit{convite} can have administrative costs from US\$500 to US\$5,200, depending on the bid complexity. Thus, if a public tender fails to obtain a supplier, it creates relevant waste for the government.


